\SourceLeaf{70}
Sit $r=2m-1$; fiet

\[
\frac{1}{2^{r-1}} = \frac{1}{2^{2m-2}},
\]

\noindent et ob

\[
\frac{1-(-1)^r}{4} = \frac{1-(-1)^{2m-1}}{4} = +\frac{1}{2},
\]

\[
{}^{r}\mathfrak{B}^{\left(\frac{1}{2}r+\frac{1-(-1)^r}{4}\right)} = {}^{2m-1}\mathfrak{B}^{\left(\frac{1}{2}(2m-1)+\frac{1}{2}\right)} = {}^{2m-1}\mathfrak{B}^{(m)},
\]

\noindent tandem

\[
\cos\frac{1-(-1)^r}{2}\varphi = \cos\frac{1-(-1)^{2m-1}}{2}\varphi = \cos\varphi.
\]

Tandem factor praefixus fit, existente $r$ numero impari, $=1$, existente pari,
$=1-\tfrac{2}{4}=\tfrac{1}{2}$, uti fas est.

Similiter, dico fore ipsius $\cos\varphi^r$ terminum indefinitum $h^{\text{tum}}$

\[
= \frac{1}{2^{r-1}}\,{}^{r}\mathfrak{B}^{\left(\frac{1}{2}r+\frac{1-(-1)^r}{4}+h\right)}\cos\left(2h+\frac{1-(-1)^r}{2}\right)\varphi.
\]

Existente enim $r$ pari $=2m$, ob

\[
\frac{1-(-1)^r}{2} = 0,
\]

\noindent fit expressio

\[
\frac{1}{2^{2m-1}}\,{}^{2m}\mathfrak{B}^{(m+h)}\cos2h\varphi;
\]

\noindent existente impari $=2m-1$, cum sit tunc

\[
\frac{1-(-1)^{2m-1}}{2} = 1,
\]

\noindent fiet eadem

\[
\frac{1}{2^{2m-2}}\,{}^{2m-1}\mathfrak{B}^{(m+h)}\cos(2h+1)\varphi,
\]

\noindent ut oportet.

Ad vitandam in scribendo molestiam pro quantitate simplici

\[
\frac{1-(-1)^r}{2}
\]

\noindent signum unicum v.\,g. $=a$, ut in analysi solet fieri, usurpari posset. Haberetur tunc,
contracta expressione
